\documentclass[../main.tex]{subfiles}

\begin{document}

\chapter{Fazit}
\vspace{-0.5cm}
\noindent\textit{\small Verfasst von Julian}
\vspace{0.5cm}

Das Teilprojekt „Naturgefahren“ im Rahmen von DeepAlpine hat erfolgreich gezeigt, dass der Einsatz von maschinellem Lernen, insbesondere durch den Random Forest Classifier, eine effektive Methode zur Identifikation und Vorhersage von Erdrutschen im alpinen Raum darstellt. \par

Mithilfe der entwickelten Modelle konnte quantitativ belegt werden, dass die Gefährdung durch gravitative Massenbewegungen infolge des Klimawandels in bestimmten Regionen deutlich zunimmt.
Der Vergleich verschiedener Zukunftsszenarien (SSP1-2.6, SSP3-7.0 und SSP5-8.5) zeigt für den Zeitraum 2081–2100 eine eindeutige Vergrößerung, aber auch eine Verschiebung der gefährdeten Flächen gegenüber dem aktuellen Referenzzeitraum (2015–2024). \par

Trotz gewisser methodischer Herausforderungen liefert der Ansatz einen entscheidenden gesellschaftlichen und wissenschaftlichen Mehrwert.
Die generierten Gefahrenkarten bilden eine essenzielle Planungsgrundlage für Gemeinden.
Sie ermöglichen es lokalen Entscheidungsträgern, Risikozonen datengestützt zu erkennen und die Sicherheit der Bevölkerung durch gezielte, vorausschauende Anpassungsmaßnahmen aufrechtzuerhalten.
Im akademischen Bereich konnte ein weiterer Beitrag über die ML-basierten Methoden zur Analyse von Naturgefahren geleistet werden.
Zudem wurde die Literatur im Bereich der Klimafolgenforschung im alpinen Raum um eine umfassende Analyse erweitert.

\end{document}